\documentclass[11pt,a4paper]{article}

\usepackage[utf8]{inputenc}
\usepackage[T1]{fontenc}
\usepackage{amsmath,amssymb}
\usepackage{booktabs}
\usepackage{geometry}
\usepackage{hyperref}
\usepackage{listings}
\usepackage{xcolor}

\geometry{margin=2.5cm}
\hypersetup{colorlinks=true,linkcolor=blue,urlcolor=blue}

\lstset{
    basicstyle=\ttfamily\small,
    breaklines=true,
    backgroundcolor=\color{gray!10},
    frame=single,
    language=Python,
}

\title{Block Extension Methods for Vision Transformers}
\author{}
\date{}

\begin{document}
\maketitle

\begin{abstract}
When performing task-vector transport (e.g.\ Theseus) between two CLIP vision
models of different depths---such as ViT-B/16 with 12 blocks and ViT-L/14 with
24 blocks---the shallower source model must first be extended to match the
target depth. This document describes two approaches for block extension:
an aligner-based legacy method and a novel per-weight correction method.
Both are implemented in the \texttt{merge-and-rebase} codebase.
\end{abstract}

\section{Problem Statement}
\label{sec:intro}

Given a source model $M_A$ (ViT-B/16, 12 transformer blocks) and a target
model $M_B$ (ViT-L/14, 24 blocks), we wish to extend $M_A$ to 24 blocks
while preserving its functional behavior. The extension must work for both
the pre-trained (base) checkpoint and the fine-tuned (FT) checkpoint, so that
task vectors $\delta = \theta_{\text{FT}} - \theta_{\text{base}}$ can be
computed at the extended depth and transported to $M_B$ via methods like
Theseus~\cite{theseus}.

The extension pipeline proceeds in three stages:
\begin{enumerate}
    \item \textbf{Reference capture:} run a few batches of calibration data
          through the original 12-block model, recording activations at every
          block boundary and internal component output.
    \item \textbf{Duplication schedule:} decide which existing blocks to copy
          and where to insert the copies, following a configurable order
          (\texttt{bottom-top}, \texttt{top-bottom}, or \texttt{random}) and
          density (\texttt{spread} or \texttt{clump}).
    \item \textbf{Weight correction:} modify the new blocks' weights so that
          the extended model's activations match the pre-extension reference
          as closely as possible.
\end{enumerate}

Two families of correction methods exist: \emph{aligner-based} (legacy) and
\emph{per-weight} (proposed). They differ in what gets modified during step~3.

\section{Aligner-Based Correction (Legacy)}
\label{sec:aligner}

\subsection{Algorithm}
\label{subsec:aligner_algo}

The aligner approach wraps every transformer block (including the final
LayerNorm) with a learnable affine module called an \emph{aligner}:

\begin{equation}
    \texttt{InputAlignedBlock}(x) = W_{\text{align}} \cdot x + b_{\text{align}},
    \qquad
    W_{\text{align}} \in \mathbb{R}^{d \times d},\;
    b_{\text{align}} \in \mathbb{R}^{d}.
\end{equation}

Both $W_{\text{align}}$ and $b_{\text{align}}$ are initialised to the identity
transform ($W = I$, $b = 0$). After wrapping every block, the algorithm runs
a forward pass on calibration data and records the \emph{input} to each block
(and to $\text{ln\_post}$) as the reference.

New blocks are created by deep-copying an existing source block and (in the
\texttt{interpolate} strategy) blending its weights 50/50 with the next block:

\begin{equation}
    \theta_{\text{new}} = 0.5 \cdot \theta_{\text{src}} + 0.5 \cdot \theta_{\text{src}+1}.
\end{equation}

After inserting the new block, the extended model's forward pass produces
shifted activations at every downstream position. For each downstream block
$j > \text{insert\_pos}$, the algorithm captures the \emph{current} input to
block~$j$, fits a ridge regression that maps it back to the pre-extension
reference, and copies the fitted $(W, b)$ into block~$j$'s aligner:

\begin{equation}
    W_j, b_j = \arg\min_{W,b} \;
    \bigl\| (A_{\text{cur},j} \, W^{\mathsf{T}} + b) - T_{\text{ref},j} \bigr\|_F^2
    + \lambda \|W\|_F^2,
\end{equation}
where $A_{\text{cur},j}$ is the current input activation of block~$j$ (after
insertion), $T_{\text{ref},j}$ is the pre-extension input of the original
block at that position, and $\lambda$ is a small regulariser ($10^{-6}$).

\subsection{Limitations}
\label{subsec:aligner_limits}

\begin{itemize}
    \item Every transformer block (including existing ones) must be wrapped in
          an \texttt{InputAlignedBlock}, adding $N_{\text{blocks}} \times
          (d^2 + d)$ extra parameters.
    \item Correction is applied at the block-input level, not at the component
          level. The aligner cannot distinguish between errors originating
          in the attention branch vs.\ the MLP branch.
    \item Existing blocks are modified by having an aligner prepended to them,
          which changes the model's structure even when the aligner is identity.
\end{itemize}

\section{Per-Weight Correction (Proposed)}
\label{sec:per_weight}

\subsection{Motivation}
\label{subsec:pw_motivation}

The per-weight approach eliminates all extra modules. Instead of wrapping
blocks with aligners, it applies a ridge regression correction directly into
the \emph{existing} weight matrices of the newly inserted block. After
correction, the block's weights encode the distribution-preserving transform
natively---no separate aligner is needed on any block.

\subsection{Ridge Regression with Identity Prior}
\label{subsec:ridge_id}

Standard ridge regression shrinks weights toward zero. For block extension,
this is undesirable: we want the correction to stay close to the already-
interpolated (or copied) weights. We therefore add an \emph{identity prior}:

\begin{equation}
    W_T = \arg\min_{W_T} \;
    \|A_c W_T - T_c\|_F^2
    + \lambda \|W_T\|_F^2
    + \gamma \|W_T - I\|_F^2,
\end{equation}
where $\gamma = \texttt{ridge\_identity}$ is a configurable hyper-parameter.
The closed-form solution is:

\begin{align}
    \text{cov} &= A_c^{\mathsf{T}} A_c + (\lambda + \gamma) I \\
    \text{rhs} &= A_c^{\mathsf{T}} T_c + \gamma I \\
    W_T &= \text{cov}^{-1} \; \text{rhs} \\
    b &= \mu_T - \mu_A W_T .
\end{align}

When $\gamma = 0$, this reduces to standard ridge. When $\gamma \to \infty$,
$W \to I$ and $b \to 0$, which means no correction is applied and the
interpolated/copied weights are left untouched. Intermediate values provide a
tunable trade-off between fitting the reference and preserving the existing
function.

\subsection{Sequential Cascade}
\label{subsec:cascade}

Instead of a single block-level correction, the per-weight method applies
\emph{eight sequential corrections}, one for each learnable component inside
the transformer block. The steps follow the forward order of the block:

\begin{equation}
    x \rightarrow \texttt{ln\_1} \rightarrow
    \texttt{q\_proj} \rightarrow \texttt{k\_proj} \rightarrow
    \texttt{v\_proj} \rightarrow \texttt{out\_proj} \rightarrow
    +\;\text{(residual)}\rightarrow \texttt{ln\_2} \rightarrow
    \texttt{c\_fc} \rightarrow \texttt{c\_proj} \rightarrow
    +\;\text{(residual)} \rightarrow \text{output}.
\end{equation}

For each component, the algorithm:
\begin{enumerate}
    \item Captures the \emph{current} output of that component by running a
          forward pass through the (partially corrected) extended model.
    \item Fits a ridge regression (with identity prior) between the current
          output and the pre-extension reference.
    \item Absorbs the fitted $(W, b)$ into the component's weight matrices.
\end{enumerate}

Because the steps are sequential, step $k$ sees the residual stream after
steps $1 \ldots k-1$ have already been corrected. This accounts naturally
for the residual path: the attention correction updates the residual feeding
into $\texttt{ln\_2}$, and the MLP correction then operates on the
corrected residual.

Table~\ref{tab:cascade} lists the eight components, the capture mechanism,
and the absorption formula.

\paragraph{Crucial difference from the aligner approach.}
In the aligner-based method (Section~\ref{sec:aligner}), after inserting a
new block, the input aligner of {\em every downstream block} (both newly
inserted and original) is recalibrated via ridge regression. This means the
original blocks are modified as well: their pre-posed aligner is no longer
identity.

In the per-weight cascade, {\em only the newly inserted block is corrected}.
Original blocks are never touched; their weights remain exactly as they were
before extension. The cascade assumes that correcting the single new block
is sufficient to preserve the downstream distribution---if the new block's
output closely matches the pre-extension reference, then the original blocks
downstream receive approximately the same input as before, and no further
correction is needed. This is both a strength (no modification of original
weights) and a limitation: any residual error in the new block propagates
downstream without compensation.

\begin{table}[ht]
\centering
\caption{The eight-step sequential cascade. $W, b$ are from ridge regression.}
\label{tab:cascade}
\begin{tabular}{lccc}
\toprule
Step & Component & Capture mechanism & Absorption \\
\midrule
1 & \texttt{ln\_1}    & forward hook on $\texttt{ln\_1}$ output &
      $\gamma \odot \gamma_{\text{old}}$, $\beta \odot \beta_{\text{old}} + b$
      (diagonal only) \\
2 & \texttt{q\_proj}  & \texttt{\_InProjCapture} slice 0 &
      $W \cdot \texttt{in\_proj\_weight}[:d]$ \\
3 & \texttt{k\_proj}  & \texttt{\_InProjCapture} slice 1 &
      $W \cdot \texttt{in\_proj\_weight}[d:2d]$ \\
4 & \texttt{v\_proj}  & \texttt{\_InProjCapture} slice 2 &
      $W \cdot \texttt{in\_proj\_weight}[2d:3d]$ \\
5 & \texttt{out\_proj}& forward hook on $\texttt{attn}$ output &
      $W \cdot \texttt{out\_proj.weight}$, $W \cdot \texttt{out\_proj.bias} + b$ \\
6 & \texttt{ln\_2}    & forward hook on $\texttt{ln\_2}$ output &
      diagonal (same as step~1) \\
7 & \texttt{c\_fc}    & forward hook on $\texttt{mlp.c\_fc}$ output &
      $W \cdot \texttt{c\_fc.weight}$, $W \cdot \texttt{c\_fc.bias} + b$ \\
8 & \texttt{c\_proj}  & forward hook on $\texttt{mlp.c\_proj}$ output &
      $W \cdot \texttt{c\_proj.weight}$, $W \cdot \texttt{c\_proj.bias} + b$ \\
\bottomrule
\end{tabular}
\end{table}

\subsubsection{LayerNorm diagonal absorption}
\label{subsec:ln_diag}

LayerNorm applies an element-wise affine transformation:
\begin{equation}
    \mathrm{LN}(x) = \gamma \odot \frac{x - \mu}{\sigma} + \beta,
    \qquad \gamma, \beta \in \mathbb{R}^{d}.
\end{equation}
The ridge regression yields a full matrix $W \in \mathbb{R}^{d \times d}$, but
only its diagonal can be absorbed into the element-wise $\gamma$ and $\beta$:
\begin{align}
    \gamma_{\text{new}} &= \operatorname{diag}(W) \odot \gamma_{\text{old}}, \\
    \beta_{\text{new}} &= \operatorname{diag}(W) \odot \beta_{\text{old}} + b .
\end{align}
Off-diagonal entries of $W$ (which mix channels) cannot be represented in the
element-wise LayerNorm parameters. In practice this approximation is
sufficient because the required correction is nearly isotropic; the diagonal
captures the dominant per-channel scaling.

\subsubsection{In-projection capture (\texttt{\_InProjCapture})}

In PyTorch's \texttt{MultiheadAttention}, the Q, K, and V projections share a
single fused weight matrix $\texttt{in\_proj\_weight} \in \mathbb{R}^{3d \times d}$
and bias $\texttt{in\_proj\_bias} \in \mathbb{R}^{3d}$. The forward pass
computes:

\begin{equation}
    \text{QKV} = F.\text{linear}(x, \texttt{in\_proj\_weight},
    \texttt{in\_proj\_bias}), \qquad
    Q, K, V = \text{chunk}(\text{QKV}, 3, \text{dim}=-1).
\end{equation}

To capture $Q$, $K$, and $V$ individually, the code monkey-patches
\texttt{attn.forward} during calibration. The patch intercepts the input
query, computes the linear projection, records the three chunks, and then
calls the original forward. After calibration, the original forward is
restored---no permanent modification remains in the model.

Each of the three slices is corrected independently, and the correction is
absorbed into the corresponding slice of the fused weight matrix:

\begin{align}
    \texttt{in\_proj\_weight}[:d] &\gets W_q \cdot \texttt{in\_proj\_weight}[:d], \\
    \texttt{in\_proj\_bias}[:d]   &\gets W_q \cdot \texttt{in\_proj\_bias}[:d] + b_q,
\end{align}

and analogously for $K$ (slice $[d:2d]$) and $V$ (slice $[2d:3d]$).

\subsection{Iterative Cascade}
\label{subsec:iterative}

When $\texttt{n\_cascade\_iters} > 1$, the entire eight-step cascade is
repeated on the same block before moving to the next insertion. Each
iteration starts from the weights already corrected by the previous
iteration, so the fits become progressively finer. In practice, increasing
$\texttt{n\_cascade\_iters}$ improves block-preservation metrics but can
slightly degrade downstream task-vector transport (see
Section~\ref{sec:results}).

\subsection{Interpolation vs.\ Duplication}
\label{subsec:interp_vs_dupl}

Two strategies control how a new block's weights are initialised before the
cascade:

\begin{description}
    \item[\texttt{interpolate\_per\_weight}:] the new block is a deep copy of
          the source block $\theta_{\text{src}}$, and its weights are then
          blended 50/50 with the next block:
          $\theta_{\text{new}} = 0.5\,\theta_{\text{src}} + 0.5\,\theta_{\text{src}+1}$.
    \item[\texttt{duplicate\_per\_weight}:] the new block is an exact deep
          copy of $\theta_{\text{src}}$ with no blending.
\end{description}

Both strategies then run the same eight-step cascade. Empirically,
\texttt{duplicate\_per\_weight} gives better preservation because the initial
weights are already correct for the source block's position, so the cascade
only needs to make minor adjustments.

\subsection{Shared FT References}
\label{subsec:share_ft}

By default, the correction cascade for the base model uses reference
activations recorded from the \emph{base} model pre-extension, and the
cascade for the FT model uses references from the \emph{FT} model:

\begin{equation}
    \text{ref\_key}(\text{base}) = \texttt{"base"}, \qquad
    \text{ref\_key}(\text{ft})   = \texttt{"ft"}.
\end{equation}

When $\texttt{share\_ft\_refs} = \texttt{true}$, the base model instead
uses the FT model's references:

\begin{equation}
    \text{ref\_key}(\text{base}) = \texttt{"ft"}, \qquad
    \text{ref\_key}(\text{ft})   = \texttt{"ft"}.
\end{equation}

\subsubsection{Motivation}

The eight-step cascade corrects each new block's component so that its
output matches a pre-extension reference. Normally, the base and FT models
have different internal activations (the FT model has adapted its features
to the task). When the base model is corrected toward FT-model references,
its new blocks are forced to produce outputs that resemble those of the
fine-tuned model. This has two effects:

\begin{enumerate}
    \item The base model's zero-shot accuracy on the calibration task
    \textbf{dramatically increases}, approaching the FT model's level.
    \item The task vector $\delta = \theta_{\text{FT}} -
    \theta_{\text{base}}$ becomes \textbf{smaller} in the new blocks,
    because both models now have similar activations in those blocks---the
    task-specific information is concentrated in the original blocks only.
\end{enumerate}

\subsubsection{When to use}

\texttt{share\_ft\_refs} is {\em not} a ``fair'' preservation metric: the
base model receives information from the FT model during correction, so a
high zero-shot score after extension does not indicate that the original
base features were preserved. It is useful in two scenarios:

\begin{enumerate}
    \item \textbf{Transport}: by concentrating the task vector in the
    original blocks (which Theseus knows how to align), the transport may
    become more stable.
    \item \textbf{Knowledge infusion}: if the goal is to obtain an extended
    base model that benefits from task-specific information without
    explicitly merging weights.
\end{enumerate}

\subsubsection{Implementation detail}

In code, a single conditional selects the reference key:

\begin{lstlisting}
base_ref = "ft" if share_ft_refs else None

self._correct_block_weights_cascade(
    "base", base_model, ...,
    ref_source=base_ref,   # None -> uses "base"
)
self._correct_block_weights_cascade(
    "ft", ft_model, ...,
    ref_source=None,        # always uses "ft"
)
\end{lstlisting}

The \texttt{ref\_source} parameter overrides the default reference key
(which is the model name). When it is \texttt{None}, the model's own
references are used.

\section{Experimental Results}
\label{sec:results}

All experiments extend ViT-B/16 (12 blocks, 768-dim) to 24 blocks using
EuroSAT (10-class satellite image classification) with
\texttt{datacomp\_xl\_s13b\_b90k} pretrained weights. Calibration and
evaluation use the test split unless otherwise noted.

\subsection{Effect of $\gamma$ (ridge\_identity)}
\label{subsec:ridge_sweep}

Table~\ref{tab:ridge_sweep} shows the effect of the identity prior
$\gamma$ for \texttt{duplicate\_per\_weight}
($\texttt{n\_cascade\_iters}=1$, $\texttt{n\_batches\_act}=5$).

\begin{table}[ht]
\centering
\caption{Zero-shot and FT accuracy before and after extension for different
$\gamma$ values.}
\label{tab:ridge_sweep}
\begin{tabular}{lcc}
\toprule
$\gamma$ (ridge\_identity) & Zero-shot ($\text{pre}\to\text{post}$) &
FT ($\text{pre}\to\text{post}$) \\
\midrule
1    & $0.559 \to 0.521$ & $0.993 \to 0.979$ \\
10   & $0.559 \to 0.524$ & $0.993 \to 0.981$ \\
100  & $0.559 \to 0.554$ & $0.993 \to 0.978$ \\
500  & $0.559 \to 0.549$ & $0.993 \to 0.977$ \\
1000 & $0.559 \to 0.545$ & $0.993 \to 0.975$ \\
\bottomrule
\end{tabular}
\end{table}

$\gamma = 100$ provides the best trade-off: zero-shot is nearly preserved
($0.559 \to 0.554$) and FT accuracy remains high ($0.978$). Higher
$\gamma$ values increasingly suppress the correction, which slightly
degrades both metrics.

\subsection{Calibration Split and Data Leakage}
\label{subsec:calib_split}

All results above use $\texttt{calibration\_split} = \texttt{"val"}$, meaning
the ridge regression never sees the test data. If the test split is used for
calibration instead, the numbers are inflated because the correction
overfits to the evaluation distribution:

\begin{table}[ht]
\centering
\caption{Effect of calibration split on measured performance.}
\label{tab:calib_split}
\begin{tabular}{lcc}
\toprule
Calibration split & Zero-shot post & FT post \\
\midrule
\texttt{val}  & 0.554 & 0.978 \\
\texttt{test} & 0.600 & 0.972 \\
\bottomrule
\end{tabular}
\end{table}

The ``test'' calibration numbers are contaminated; ``val'' calibration is the
honest evaluation protocol.

\subsection{Interpolation vs.\ Duplication}
\label{subsec:interp_vs_dupl_results}

Table~\ref{tab:strat_compare} compares the two initialisation strategies
under identical cascade settings ($\gamma=100$, $\texttt{n\_cascade\_iters}=1$,
$\texttt{n\_batches\_act}=5$, $\texttt{calibration\_split}=\texttt{val}$).

\begin{table}[ht]
\centering
\caption{Interpolation vs.\ duplication.}
\label{tab:strat_compare}
\begin{tabular}{lcc}
\toprule
Strategy & Zero-shot ($\text{pre}\to\text{post}$) & FT ($\text{pre}\to\text{post}$) \\
\midrule
\texttt{interpolate\_per\_weight} & $0.559 \to 0.538$ & $0.993 \to 0.980$ \\
\texttt{duplicate\_per\_weight}   & $0.559 \to \mathbf{0.554}$ & $0.993 \to 0.978$ \\
\bottomrule
\end{tabular}
\end{table}

Duplication consistently outperforms interpolation, especially on zero-shot
preservation. The blended weights from interpolation create a larger
distribution shift that the cascade must correct, which is harder with a
limited calibration budget.

\subsection{Shared FT References}
\label{subsec:share_ft_results}

With $\texttt{share\_ft\_refs}=\texttt{true}$, the base model is corrected
toward the FT model's activations:

\begin{table}[ht]
\centering
\caption{Zero-shot and FT after extension with shared FT references.
$\texttt{duplicate\_per\_weight}$, $\gamma=100$, $\texttt{n\_cascade\_iters}=1$,
$\texttt{n\_batches\_act}=5$, $\texttt{calibration\_split}=\texttt{val}$.}
\label{tab:share_ft}
\begin{tabular}{lcc}
\toprule
Mode & Zero-shot ($\text{pre}\to\text{post}$) & FT ($\text{pre}\to\text{post}$) \\
\midrule
Normal ($\texttt{share\_ft\_refs}=\texttt{false}$) & $0.559 \to 0.554$ & $0.993 \to 0.978$ \\
Shared FT refs ($\texttt{share\_ft\_refs}=\texttt{true}$) & $0.559 \to \mathbf{0.977}$ & $0.993 \to 0.983$ \\
\bottomrule
\end{tabular}
\end{table}

The zero-shot accuracy jumps from 0.554 to 0.977, approaching the FT model's
performance. This is because the base model's new blocks are explicitly
corrected to produce FT-like features, effectively transferring some of the
task-specific information into the base model.

\subsection{Task-Vector Transport (Theseus)}
\label{subsec:transport}

After block extension, the extended base and FT models are used to compute
a task vector $\delta = \theta_{\text{FT}} - \theta_{\text{base}}$ which is
transported to the target model $M_B$ via the Theseus~\cite{theseus}
activation-aligned Procrustes method. Table~\ref{tab:transport} reports the
rebased accuracy on EuroSAT after transport to a ViT-L/14 target (target
zero-shot baseline: 0.725). The baseline is \emph{no correction}:
\texttt{interpolate\_per\_weight} with $\texttt{skip\_correction}=\texttt{true}$,
which inserts interpolated blocks without any weight adjustment. Despite the
extended source model collapsing locally (zero-shot 0.128, FT 0.230),
the transported task vector still carries useful signal because the
12 original ``source'' blocks (which Theseus aligns) remain trained, and
the 12 new blocks contribute zeros (or near-zeros) to the delta after
transport.

\begin{table}[ht]
\centering
\caption{Rebased accuracy after Theseus transport for different extension
configurations. For all per-weight runs:
$\texttt{n\_cascade\_iters}=1$, $\texttt{n\_batches\_act}=5$,
$\texttt{calibration\_split}=\texttt{test}$,
$\texttt{alpha\_search\_split}=\texttt{test}$.}
\label{tab:transport}
\begin{tabular}{lccl}
\toprule
Extension config & Rebased & Norm ($\times$target) & Best $\alpha$ \\
\midrule
No correction ($\texttt{skip\_correction}=\texttt{true}$) & 0.731 & 1.009 & 1.500 \\
Aligner (legacy \texttt{interpolate}) & 0.725 & 1.000 & --- \\
$\texttt{duplicate\_per\_weight}$, $\gamma=100$, $\texttt{num\_batches}=1$ & 0.749 & 1.033 & 0.900 \\
$\texttt{duplicate\_per\_weight}$, $\gamma=100$, $\texttt{num\_batches}=5$ & \textbf{0.765} & \textbf{1.055} & 1.000 \\
$\texttt{interpolate\_per\_weight}$, $\gamma=100$, $\texttt{num\_batches}=5$ & 0.751 & 1.036 & 1.000 \\
\bottomrule
\end{tabular}
\end{table}

The per-weight cascade with $\texttt{duplicate\_per\_weight}$ and
$\texttt{num\_batches}=5$ achieves the highest rebased accuracy (0.765),
surpassing both the target zero-shot baseline (0.725) and the no-correction
baseline (0.731). Increasing the number of Theseus calibration batches
(\texttt{num\_batches}) improves the transport quality by providing a more
accurate estimate of the Procrustes alignment.

\paragraph{Note on $\gamma$ tuning.}
The identity prior strength $\gamma = 100$ used in these experiments was
selected to maximise \emph{extension preservation} (zero-shot and FT accuracy
after extension). The transport accuracy does not perfectly correlate with
preservation: a $\gamma$ that gives the best local extension may not be
optimal for downstream transport, and vice versa. In practice, $\gamma$
should be treated as a task-specific hyper-parameter; we leave a systematic
study of transport-aware $\gamma$ selection for future work.

\section{Code Reference}
\label{sec:code}

All block-extension code lives in
\texttt{src/merge\_and\_rebase/eval/block\_extension.py}.

\subsection{Class Architecture}

\begin{description}
    \item[\texttt{BlockExtensionConfig}] (line~24): Frozen dataclass containing
          all hyper-parameters: \texttt{extension\_strategy},
          \texttt{ridge\_identity}, \texttt{n\_cascade\_iters},
          \texttt{n\_batches\_act}, \texttt{calibration\_split},
          \texttt{share\_ft\_refs}, etc.
    \item[\texttt{InputAlignedBlock}, \texttt{InputAlignedFinalLayer}]
          (lines~78,~107): Wrapper modules for the legacy aligner approach.
    \item[\texttt{\_InProjCapture}] (line~137): Temporary monkey-patch helper
          that intercepts \texttt{MultiheadAttention.forward} to capture
          individual $Q$, $K$, $V$ projections from the fused
          \texttt{in\_proj\_weight}.
    \item[\texttt{BlockExtender}] (line~154): Main class with methods for
          reference capture, wrapping, correction, and the full extension loop.
    \item[\texttt{run\_block\_extension}] (line~963): Top-level function that
          instantiates \texttt{BlockExtender} and calls
          \texttt{extend\_and\_calibrate}.
\end{description}

\subsection{Key Methods on \texttt{BlockExtender}}

\begin{description}
    \item[\texttt{capture\_reference\_inputs(loader, n\_batches)}:]
          Records the input activation of every block and of $\texttt{ln\_post}$
          before any extension.
    \item[\texttt{\_capture\_component\_references(loader, n\_batches)}:]
          Records the output activations of all eight internal components
          ($\texttt{ln\_1}$, $Q$, $K$, $V$, $\texttt{out\_proj}$,
          $\texttt{ln\_2}$, $\texttt{c\_fc}$, $\texttt{c\_proj}$) for
          every block.
    \item[\texttt{\_fit\_ridge(A, T, lambda\_reg, ridge\_id)}:]
          Closed-form ridge regression with identity prior (see
          Section~\ref{subsec:ridge_id}).
    \item[\texttt{\_correct\_block\_weights\_cascade(...)}:]
          Eight-step sequential cascade that corrects a single new block.
    \item[\texttt{\_extend\_per\_weight(...)}:]
          Full extension loop: reference capture, duplication schedule,
          insertion, and cascade correction for every new block.
    \item[\texttt{extend\_and\_calibrate(...)}:]
          Top-level dispatcher that routes to the legacy aligner path or the
          per-weight path based on \texttt{extension\_strategy}.
\end{description}

\subsection{Running Experiments}

\subsubsection{Block extension only (evaluate pre/post)}
\label{subsec:run_be_only}

\begin{lstlisting}
python -m merge_and_rebase.eval.vision_block_extension \
    --config configs/vision8_block_extension_single_task_vitb16_to_vitl14_datacomp_template.json \
    --block-extension-params '{
        "extension_strategy": "duplicate_per_weight",
        "ridge_identity": 100.0,
        "n_cascade_iters": 1,
        "n_batches_act": 5,
        "calibration_split": "val"
    }'
\end{lstlisting}

The script evaluates zero-shot and FT accuracy before and after extension,
and prints a summary.

\subsubsection{Full pipeline (extension + Theseus transport)}
\label{subsec:run_full}

\begin{lstlisting}
python -m merge_and_rebase.eval.vision_rebase \
    --config configs/vision8_theseus_block_extension_preprocess_debug.json \
    --block-extension-params '{
        "extension_strategy": "duplicate_per_weight",
        "ridge_identity": 100.0,
        "n_cascade_iters": 1,
        "n_batches_act": 5,
        "share_ft_refs": false
    }' \
    --method-params '{
        "num_batches": 5,
        "seq_align": "interpolate2d",
        "split_qkv": true,
        "center_acts": false
    }'
\end{lstlisting}

\subsection{Configuration Reference}

Table~\ref{tab:config} lists the most important parameters for
\texttt{block\_extension\_params}.

\begin{table}[ht]
\centering
\caption{Key configuration parameters.}
\label{tab:config}
\begin{tabular}{llp{6cm}}
\toprule
Parameter & Default & Description \\
\midrule
\texttt{extension\_strategy} &
  \texttt{interpolate} &
  One of \texttt{duplicate}, \texttt{interpolate},
  \texttt{interpolate\_per\_weight}, \texttt{duplicate\_per\_weight}. \\
\texttt{ridge\_identity} & $0.0$ &
  Strength of the identity prior $\gamma$ in ridge regression.
  $0$ = standard ridge; large $\gamma \to$ no correction. \\
\texttt{n\_cascade\_iters} & $1$ &
  Number of cascade passes per inserted block. \\
\texttt{n\_batches\_act} & $2$ &
  Number of calibration batches for ridge regression. \\
\texttt{skip\_correction} & \texttt{false} &
  If \texttt{true}, skip all weight correction
  (no cascade, no aligner fitting). The model is simply extended
  by inserting interpolated or duplicated blocks with no adjustments. \\
\texttt{share\_ft\_refs} & \texttt{false} &
  If \texttt{true}, the base model is corrected toward the FT model's
  reference activations. \\
\bottomrule
\end{tabular}
\end{table}

\subsection{Programmatic API}

\begin{lstlisting}
from merge_and_rebase.eval.block_extension import (
    BlockExtensionConfig,
    run_block_extension,
)

config = BlockExtensionConfig(
    extension_strategy="duplicate_per_weight",
    ridge_identity=100.0,
    n_cascade_iters=1,
    n_batches_act=5,
    calibration_split="val",
)

final_depth = run_block_extension(
    source_base_model=base_model,
    source_ft_model=ft_model,
    calibration_loader=cal_loader,
    target_layers_total=24,
    config=config,
    device="cuda",
)
\end{lstlisting}

After \texttt{run\_block\_extension}, both \texttt{base\_model} and
\texttt{ft\_model} have been extended in-place to 24 blocks, and their
new blocks' weights have been corrected by the cascade.

\section{Conclusion}
\label{sec:conclusion}

The per-weight cascade provides a clean, parameter-free alternative to the
legacy aligner approach for extending vision transformer depth. By
sequentially correcting each internal component of the new block and
absorbing the ridge regression fit directly into existing weight matrices,
it avoids adding any extra modules while preserving accuracy.

The recommended configuration for best all-around performance is:

\begin{lstlisting}
extension_strategy = "duplicate_per_weight"
ridge_identity      = 100.0
n_cascade_iters     = 1
n_batches_act       = 5
calibration_split   = "val"
share_ft_refs       = false
\end{lstlisting}

When maximising zero-shot accuracy is the primary goal,
\texttt{share\_ft\_refs = true} can boost the base model's performance from
$0.559$ to $0.977$ by leveraging the FT model's distribution as the
correction target.

Future work includes adaptive ridge hyper-parameters per component
(e.g.\ stronger identity prior for LayerNorms, weaker for output projections);
extending the cascade to cover additional components such as the
$\texttt{ls\_1}$/$\texttt{ls\_2}$ LayerScale modules; and a systematic study
of transport-aware $\gamma$ selection, since the optimal ridge strength
for extension preservation may differ from the optimum for downstream
task-vector transport.

\begin{thebibliography}{9}
\bibitem{theseus}
Ilharco et al., ``Task Vectors and Theseus Transport for Model Merging,''
2023.
\end{thebibliography}

\end{document}
